% !TEX root = ../../ctfp-print.tex

\lettrine[lhang=0.17]{A}{t the risk of sounding} 像个坏掉的唱片般重复，我还是要说：函子（functor）是一个非常简单却强大的概念。范畴论中充满了这类简单而有力的思想。函子是范畴之间的映射。给定两个范畴 $\cat{C}$ 和 $\cat{D}$，一个函子 $F$ 将 $\cat{C}$ 中的对象映射到 $\cat{D}$ 中的对象——这是对象上的函数。如果 $a$ 是 $\cat{C}$ 中的对象，我们将其在 $\cat{D}$ 中的像写作 $F a$（不带括号）。但范畴不仅仅是对象——它由对象和连接它们的态射组成。函子也映射态射——它是态射上的函数。但它不会随意映射态射——它会保持连接关系。因此若 $\cat{C}$ 中有一个态射 $f$ 连接对象 $a$ 到对象 $b$，
$$f \Colon a \to b$$
那么 $f$ 在 $\cat{D}$ 中的像 $F f$ 也会连接 $a$ 的像到 $b$ 的像：
$$F f \Colon F a \to F b$$

（这混合了数学符号和Haskell记法，希望你现在已熟悉这种表示方式。当对对象或态射应用函子时，我不会使用括号。）

\begin{figure}[H]
  \centering\includegraphics[width=0.3\textwidth]{images/functor.jpg}
\end{figure}

\noindent
如图所示，函子保持了范畴的结构：在一个范畴中相连的对象在另一个范畴中仍然相连。但范畴的结构还有更多要素：态射的复合关系。若 $h$ 是 $f$ 和 $g$ 的复合：
$$h = g \circ f$$
我们要求其在 $F$ 下的像也是 $f$ 和 $g$ 像的复合：
$$F h = F g \circ F f$$

\begin{figure}[H]
  \centering
  \includegraphics[width=0.3\textwidth]{images/functorcompos.jpg}
\end{figure}

\noindent
最后，我们要求 $\cat{C}$ 中的所有恒等态射都被映射到 $\cat{D}$ 中的恒等态射：
$$F \idarrow[a] = \idarrow[F a]$$

\noindent
这里，$\idarrow[a]$ 是对象 $a$ 上的恒等态射，$\idarrow[F a]$ 是 $F a$ 上的恒等态射。

\begin{figure}[H]
  \centering
  \includegraphics[width=0.3\textwidth]{images/functorid.jpg}
\end{figure}

\noindent
请注意，这些条件使得函子比普通函数更加受限。函子必须保持范畴的结构。如果你将范畴想象为由态射网络连接的一组对象，函子不允许在这个结构上引入任何撕裂。它可以将对象压合在一起，可以将多个态射粘合成一个，但永远不会将事物分开。这种"无撕裂"约束类似于你在微积分中熟知的连续性条件。从这个意义上说，函子是"连续的"（尽管存在更严格的连续性定义）。就像函数一样，函子既可以做坍缩也可以做嵌入。当源范畴远小于目标范畴时，嵌入特性更为明显。极端情况下，源范畴可以是平凡的单元素范畴——只有一个对象和一个态射（恒等态射）。从单元素范畴到其他范畴的函子仅仅选择该范畴中的一个对象。这完全类似于从单元素集合出发的态射选择目标集合中的元素。最大化的坍缩函子称为常值函子 $\Delta_c$。它将源范畴中的每个对象映射到目标范畴中选定的对象 $c$，也将每个态射映射到恒等态射 $\idarrow[c]$。它像黑洞一样将一切压缩成一个奇点。当我们讨论极限和余极限时会再次见到这个函子。

\section{编程中的函子}

让我们回归现实谈谈编程。我们有类型和函数构成的范畴。我们可以讨论将此范畴映射到自身的函子——这类函子称为自函子（endofunctor）。那么在类型范畴中什么是自函子？首先，它将类型映射到类型。我们见过这类映射的例子，可能当时并未意识到它们就是自函子。我说的是那些由其他类型参数化的类型定义。来看几个例子。

\subsection{Maybe 函子}

\code{Maybe} 的定义是从类型 \code{a} 到类型 \code{Maybe a} 的映射：

\src{snippet01}
一个重要细节：\code{Maybe} 本身不是类型，它是\emph{类型构造器}。你需要给它提供一个类型参数（如 \code{Int} 或 \code{Bool}）才能生成具体类型。\code{Maybe} 不带参数时表示类型上的函数。但我们能将 \code{Maybe} 变成函子吗？（从现在起，当我在编程上下文中谈论函子时，几乎总是指自函子。）函子不仅是对象（此处是类型）的映射，还是态射（此处是函数）的映射。对于任意从 \code{a} 到 \code{b} 的函数：

\src{snippet02}
我们想生成一个从 \code{Maybe a} 到 \code{Maybe b} 的函数。定义这样的函数需要考虑两种情况，对应 \code{Maybe} 的两个构造器。\code{Nothing} 情况很简单：直接返回 \code{Nothing}。若参数是 \code{Just}，我们对其内容应用函数 \code{f}。因此 \code{f} 在 \code{Maybe} 下的像是函数：

\src{snippet03}
（顺便提一下，在Haskell中变量名可以包含撇号，这种情况特别方便。）在Haskell中，我们将函子的态射映射部分实现为名为 \code{fmap} 的高阶函数。对于 \code{Maybe}，它的签名如下：

\src{snippet04}

\begin{figure}[H]
  \centering
  \includegraphics[width=0.35\textwidth]{images/functormaybe.jpg}
\end{figure}

\noindent
我们常说 \code{fmap} \emph{提升}了一个函数。提升后的函数作用于 \code{Maybe} 值。通常由于柯里化，这个签名可以有两种解读：作为接受一个函数 \code{(a -> b)} 并返回函数 \code{(Maybe a -> Maybe b)} 的一元函数；或者作为接受两个参数并返回 \code{Maybe b} 的二元函数：

\src{snippet05}
基于之前的讨论，这是我们为 \code{Maybe} 实现的 \code{fmap}：

\src{snippet06}
要证明类型构造器 \code{Maybe} 与函数 \code{fmap} 共同构成函子，我们需要证明 \code{fmap} 保持恒等和复合关系。这些被称为"函子定律"，它们本质上确保了范畴结构的保持。

\subsection{方程推理}

为了证明函子定律，我将使用\newterm{方程推理}，这是Haskell中常见的证明技巧。它利用了Haskell函数定义为等式的特点：左侧等于右侧。你可以随时用一个替换另一个，必要时重命名变量以避免名称冲突。可以视为内联函数，或者反向重构表达式为函数。以恒等函数为例：

\src{snippet07}
例如，若在某个表达式中看到 \code{id y}，你可以将其替换为 \code{y}（内联）。进一步地，若看到 \code{id} 应用于表达式如 \code{id (y + 2)}，你可以将其替换为表达式本身 \code{(y + 2)}。这种替换是双向的：你可以用 \code{id e} 替换任何表达式 \code{e}（重构）。若函数通过模式匹配定义，你可以独立使用每个子定义。例如，基于上述 \code{fmap} 定义，你可以将 \code{fmap f Nothing} 替换为 \code{Nothing}，反之亦然。让我们看实际效果。先从保持恒等开始：

\src{snippet08}
需要考虑两种情况：\code{Nothing} 和 \code{Just}。第一种情况如下（我使用Haskell伪代码将左侧转换为右侧）：

\begin{snip}{haskell}
  fmap id Nothing
= { fmap 的定义 }
  Nothing
= { id 的定义 }
  id Nothing
\end{snip}
注意最后一步我逆向使用了 \code{id} 的定义。将表达式 \code{Nothing} 替换为 \code{id\ Nothing}。实践中，这种证明通过"两头烧蜡烛"的方式进行，直到中间得到相同表达式——这里是 \code{Nothing}。第二种情况同样简单：

\begin{snip}{haskell}
  fmap id (Just x)
= { fmap 的定义 }
  Just (id x)
= { id 的定义 }
  Just x
= { id 的定义 }
  id (Just x)
\end{snip}
现在证明 \code{fmap} 保持复合关系：

\src{snippet09}
先处理 \code{Nothing} 情况：

\begin{snip}{haskell}
  fmap (g . f) Nothing
= { fmap 的定义 }
  Nothing
= { fmap 的定义 }
  fmap g Nothing
= { fmap 的定义 }
  fmap g (fmap f Nothing)
\end{snip}
再处理 \code{Just} 情况：

\begin{snip}{haskell}
  fmap (g . f) (Just x)
= { fmap 的定义 }
  Just ((g . f) x)
= { 复合的定义 }
  Just (g (f x))
= { fmap 的定义 }
  fmap g (Just (f x))
= { fmap 的定义 }
  fmap g (fmap f (Just x))
= { 复合的定义 }
  (fmap g . fmap f) (Just x)
\end{snip}
值得强调的是，方程推理对C++风格的副作用函数无效。考虑以下代码：

\begin{snip}{cpp}
int square(int x) {
    return x * x;
}

int counter() {
    static int c = 0;
    return c++;
}

double y = square(counter());
\end{snip}
使用方程推理，你可以内联 \code{square} 得到：

\begin{snip}{cpp}
double y = counter() * counter();
\end{snip}
这显然不是有效的转换，结果也不会相同。尽管如此，若将 \code{square} 实现为宏，C++编译器仍会尝试使用方程推理，导致灾难性后果。

\subsection{Optional}

函子在Haskell中易于表达，但在任何支持泛型编程和高阶函数的语言中都可以定义。让我们考虑 \code{Maybe} 的C++对应物——模板类型 \code{optional}。以下是其实现草图（实际实现复杂得多，需处理各种传参方式、拷贝语义及C++特有的资源管理问题）：

\begin{snip}{cpp}
template<class T>
class optional {
    bool _isValid; // 标记位
    T _v;
public:
    optional()    : _isValid(false) {}        // Nothing
    optional(T x) : _isValid(true) , _v(x) {} // Just
    bool isValid() const { return _isValid; }
    T val() const { return _v; } };
\end{snip}
这个模板提供了函子定义的一部分：类型的映射。它将任何类型 \code{T} 映射到新类型 \code{optional<T>}。我们定义其对函数的作用：

\begin{snip}{cpp}
template<class A, class B>
std::function<optional<B>(optional<A>)>
fmap(std::function<B(A)> f) {
    return [f](optional<A> opt) {
        if (!opt.isValid())
            return optional<B>{};
        else
            return optional<B>{ f(opt.val()) };
    };
}
\end{snip}
这是一个高阶函数，接受函数作为参数并返回函数。以下是其非柯里化版本：

\begin{snip}{cpp}
template<class A, class B>
optional<B> fmap(std::function<B(A)> f, optional<A> opt) {
    if (!opt.isValid())
        return optional<B>{};
    else
        return optional<B>{ f(opt.val()) };
}
\end{snip}
还可以让 \code{fmap} 成为 \code{optional} 的模板方法。这种多重选择使得在C++中抽象函子模式变得困难。函子应该是一个继承接口（遗憾的是无法拥有模板虚函数）？应该是柯里化还是非柯里化的自由模板函数？C++编译器能否正确推断缺失类型，还是需要显式指定？假设输入函数 \code{f} 接受 \code{int} 返回 \code{bool}，编译器如何确定 \code{g} 的类型：

\begin{snip}{cpp}
auto g = fmap(f);
\end{snip}
尤其当未来存在多个重载 \code{fmap} 的函子时？（我们很快会见到更多函子。）

\subsection{类型类}

Haskell如何抽象函子呢？它使用类型类机制。类型类定义支持共同接口的一组类型。例如，支持相等性的类定义如下：

\src{snippet10}
此定义表明，类型 \code{a} 若支持运算符 \code{(==)}（接受两个 \code{a} 参数并返回 \code{Bool}），则属于 \code{Eq} 类。若要告诉Haskell某个特定类型是 \code{Eq}，需要声明其为此类的\newterm{实例}并提供 \code{(==)} 实现。例如，给定二维 \code{Point} 的定义（两个 \code{Float} 的乘积类型）：

\src{snippet11}
可以定义点的相等性：

\src{snippet12}
这里我在两个模式 \code{(Pt x y)} 和 \code{(Pt x' y')} 之间使用了运算符 \code{(==)}（即正在定义的运算符）。函数体跟在单个等号之后。一旦 \code{Point} 被声明为 \code{Eq} 实例，就可以直接比较点的相等性。注意，与C++或Java不同，定义 \code{Point} 时无需指定 \code{Eq} 类（或接口）——可以在客户端代码中稍后指定。类型类也是Haskell唯一的函数（和运算符）重载机制。我们需要用它来为不同函子重载 \code{fmap}。不过有个复杂之处：函子不是作为类型定义的，而是作为类型的映射（类型构造器）。我们需要一个类型类，不像 \code{Eq} 那样是类型族，而是类型构造器族。幸运的是，Haskell类型类既适用于类型构造器也适用于类型。以下是 \code{Functor} 类的定义：

\src{snippet13}
它规定，若存在具有指定类型签名的函数 \code{fmap}，则 \code{f} 是 \code{Functor}。小写 \code{f} 是类型变量，类似类型变量 \code{a} 和 \code{b}。编译器可通过其用法（如 \code{f a} 和 \code{f b} 中对其他类型的施用）推断出它代表类型构造器而非类型。相应地，声明 \code{Functor} 实例时，需要给出类型构造器，如 \code{Maybe}：

\src{snippet14}
顺便说，\code{Functor} 类及其许多简单数据类型的实例定义（包括 \code{Maybe}）都是标准预导入库的一部分。

\subsection{C++中的函子}

我们能在C++中采用相同方法吗？类型构造器对应模板类如 \code{optional}，因此类比之下，我们需要用\newterm{模板模板参数} \code{F} 参数化 \code{fmap}。其语法如下：

\begin{snip}{cpp}
template<template<class> F, class A, class B>
F<B> fmap(std::function<B(A)>, F<A>);
\end{snip}
我们希望为不同函子特化此模板。不幸的是，C++禁止部分特化模板函数。你不能写：

\begin{snip}{cpp}
template<class A, class B>
optional<B> fmap<optional>(std::function<B(A)> f, optional<A> opt)
\end{snip}
相反，我们必须依赖函数重载，这使我们回到原始的非柯里化 \code{fmap} 定义：

\begin{snip}{cpp}
template<class A, class B>
optional<B> fmap(std::function<B(A)> f, optional<A> opt) {
    if (!opt.isValid())
        return optional<B>{};
    else
        return optional<B>{ f(opt.val()) };
}
\end{snip}
这个定义有效，但仅仅因为 \code{fmap} 的第二个参数选择了重载。它完全忽略了更通用的 \code{fmap} 定义。

\subsection{列表函子}

为了直观理解函子在编程中的作用，我们需要更多例子。任何由其他类型参数化的类型都是函子的候选者。通用容器由存储元素的类型参数化，因此让我们看一个简单的容器：列表：

\src{snippet15}
我们有类型构造器 \code{List}，它将任何类型 \code{a} 映射到类型 \code{List a}。要证明 \code{List} 是函子，需要定义函数的提升：给定函数 \code{a -> b}，定义函数 \code{List a -> List b}：

\src{snippet16}
作用于 \code{List a} 的函数必须考虑对应两个列表构造器的两种情况。\code{Nil} 情况很直接——只需返回 \code{Nil}。空列表没什么可做的。\code{Cons} 情况较复杂，因为它涉及递归。让我们暂时退一步思考我们在做什么。我们有一个 \code{a} 的列表，一个将 \code{a} 转换为 \code{b} 的函数 \code{f}，想要生成一个 \code{b} 的列表。显然的做法是用 \code{f} 将列表的每个元素从 \code{a} 转换为 \code{b}。如何在实践中实现这一点？给定（非空）列表定义为头部和尾部的 \code{Cons}，我们对头部应用 \code{f}，并对尾部应用提升后的（\code{fmap} 处理的）\code{f}。这是一个递归定义，因为我们正在用提升后的 \code{f} 定义提升后的 \code{f}：

\src{snippet17}
注意，在右侧，\code{fmap f} 应用于比当前定义的列表更短的列表——应用于其尾部。我们递归处理越来越短的列表，最终必然到达空列表 \code{Nil}。正如之前决定的，\code{fmap f} 作用于 \code{Nil} 返回 \code{Nil}，从而终止递归。要得到最终结果，我们用 \code{Cons} 构造器组合新头部 \code{(f x)} 和新尾部 \code{(fmap f t)}。综合起来，以下是列表函子的实例声明：

\src{snippet18}
若你更熟悉C++，考虑 \code{std::vector} 的情况，它可能是最通用的C++容器。\code{std::vector} 的 \code{fmap} 实现只是 \code{std::transform} 的薄封装：

\begin{snip}{cpp}
template<class A, class B>
std::vector<B> fmap(std::function<B(A)> f, std::vector<A> v) {
    std::vector<B> w;
    std::transform( std::begin(v)
                  , std::end(v)
                  , std::back_inserter(w)
                  , f);
    return w;
}
\end{snip}
例如，可以用它平方数字序列的元素：

\begin{snip}{cpp}
std::vector<int> v{ 1, 2, 3, 4 };
auto w = fmap([](int i) { return i*i; }, v);
std::copy( std::begin(w)
         , std::end(w)
         , std::ostream_iterator(std::cout, ", "));
\end{snip}
大多数C++容器通过实现可传递给 \code{std::transform} 的迭代器成为函子，\code{std::transform} 是 \code{fmap} 更原始的表亲。不幸的是，函子的简洁性通常淹没在迭代器和临时对象的杂乱代码中（见上面 \code{fmap} 的实现）。很高兴地说，新的C++提案范围库使范围的函子性质更加突出。

\subsection{读者函子}

现在你可能已经形成了一些直觉——例如，函子像某种容器——让我展示一个乍看截然不同的例子。考虑将类型 \code{a} 映射到返回 \code{a} 的函数类型的映射。我们尚未深入讨论函数类型——完整的范畴论处理将在后面——但作为程序员我们对其有一定理解。在Haskell中，函数类型通过箭头类型构造器 \code{(->)} 构建，它接受两个类型：参数类型和结果类型。你已经见过它的中缀形式 \code{a -> b}，但同样可以前缀使用：

\src{snippet19}
就像常规函数一样，多参数类型函数可以部分应用。因此当我们只为箭头提供一个类型参数时，它仍期待另一个参数。这就是为什么：

\src{snippet20}
是一个类型构造器。它还需要一个类型 \code{b} 才能生成完整类型 \code{a -> b}。目前，它定义了一族由 \code{a} 参数化的类型构造器。让我们看看这是否也是函子族。处理两个类型参数可能会有些混乱，所以让我们进行一些重命名。将参数类型命名为 \code{r}，结果类型命名为 \code{a}，与之前的函子定义一致。因此我们的类型构造器接受任何类型 \code{a} 并将其映射到类型 \code{r -> a}。要证明它是函子，我们需要将函数 \code{a -> b} 提升为接受 \code{r -> a} 并返回 \code{r -> b} 的函数。这些类型分别通过类型构造器 \code{(->) r} 对 \code{a} 和 \code{b} 施用得到。以下是此情况下的 \code{fmap} 类型签名：

\src{snippet21}
我们需要解决以下难题：给定函数 \code{f :: a -> b} 和函数 \code{g :: r -> a}，创建函数 \code{r -> b}。我们只能有一种方式组合这两个函数，结果正好是我们需要的。以下是我们的 \code{fmap} 实现：

\src{snippet22}
它完美工作！如果你喜欢简洁的记法，可以通过注意到组合可以前缀形式重写进一步简化定义：

\src{snippet23}
省略参数可以得到两个函数的直接等式：

\src{snippet24}
这种类型构造器 \code{(->) r} 与上述 \code{fmap} 实现的组合称为读者函子（reader functor）。

\section{作为容器的函子}

我们已经见过编程语言中的一些函子示例，它们定义了通用目的的容器，或者至少是包含其参数化类型的值的对象。读者函子似乎是个例外，因为我们通常不认为函数是数据。但我们已经看到纯函数可以被记忆化，而函数执行可以转化为查表操作。表格就是数据。相反地，由于Haskell的惰性求值特性，一个传统的容器如列表实际上可能被实现为函数。例如，自然数的无限列表就可以紧凑地定义为：

\src{snippet25}
在第一行中，方括号是Haskell内置的列表类型构造器。第二行中方括号用于创建列表字面量。显然这样的无限列表无法存储在内存中。编译器将其实现为按需生成\code{Integer}的函数。Haskell有效地模糊了数据与代码的界限。列表可以被视为函数，而函数又可以被视为将参数映射到结果的表格。当函数的定义域是有限且规模不大时，这种实现方式甚至可能是实用的。然而将\code{strlen}实现为查表法就不可行，因为存在无限多不同的字符串。作为程序员，我们不喜欢处理无穷大，但在范畴论中你得学会习惯与无穷大共存。无论是所有字符串的集合还是包含宇宙过去、现在和未来所有可能状态的集合——我们都能处理！因此我喜欢将函子对象（由自函子生成的类型对象）视为包含其参数化类型的值或值集合的容器，即使这些值并非物理上真实存在。一个函子示例是C++的\code{std::future}，它可能在某个时刻包含某个值（但不保证），访问该值时可能需要阻塞等待其他线程完成执行。另一个例子是Haskell的\code{IO}对象，它可能包含用户输入，或显示着"Hello World!"的未来宇宙状态。根据这种诠释，函子对象是可能包含其参数化类型的值或值生成方案的实体。我们完全不关心能否实际访问这些值——这是可选功能，超出函子范畴。我们唯一感兴趣的是能够通过函数来操作这些值。若值可访问，则应能观察到操作结果；若不可访问，则只需关注操作的组合律成立，且恒等函数的操作不会改变任何东西。为了展示我们对访问函子对象内部值的漠视程度，请看这个完全忽略其参数\code{a}的类型构造器：

\src{snippet26}
\code{Const}类型构造器接受两个类型\code{c}和\code{a}。就像我们对箭头构造器所做的那样，我们将部分应用它来创建函子。数据构造器（同样称为\code{Const}）仅接受一个\code{c}类型的值，它与\code{a}完全无关。该类型构造器的\code{fmap}类型为：

\src{snippet27}
由于函子忽略其类型参数，\code{fmap}的实现可以自由忽略其函数参数——函数无物可作用：

\src{snippet28}
在C++中这个实现可能更清晰些（我从未想过自己会说出这样的话！），那里类型参数（编译期）和值（运行期）有更强区分：

\begin{snip}{cpp}
template<class C, class A>
struct Const {
    Const(C v) : _v(v) {}
    C _v;
};
\end{snip}
\code{fmap}的C++实现同样忽略函数参数，本质上直接重新封装\code{Const}参数而不改变其值：

\begin{snip}{cpp}
template<class C, class A, class B>
Const<C, B> fmap(std::function<B(A)> f, Const<C, A> c) {
    return Const<C, B>{c._v};
}
\end{snip}
尽管看似怪异，\code{Const}函子在许多构造中扮演重要角色。在范畴论中，它是前面提到的$\Delta_c$函子的特例——黑洞般的自函子情形。我们将在后续章节中再次见到它。

\section{函子组合}

不难说服自己，范畴间的函子可以像集合间的函数一样进行组合。两个函子的组合作用于对象时，就是它们各自对象映射的组合；作用于态射时同理。经过两个函子的传递后，恒等态射仍保持为恒等态射，态射的复合也保持为复合态射。实际上这个过程非常直接。特别地，自函子（endofunctor）的组合尤其简单。还记得\code{maybeTail}函数吗？我将用Haskell内置的列表实现重写它：

\src{snippet29}
（我们曾称为\code{Nil}的空列表构造器被替换为空方括号\code{{[}{]}}，\code{Cons}构造器被替换为中缀操作符\code{:}（冒号）。）\code{maybeTail}的结果类型是两个函子\code{Maybe}和\code{{[}{]}}作用于\code{a}的组合类型。每个函子都有各自的\code{fmap}实现，但如果我们想对复合结构（即Maybe列表）中的元素应用某个函数\code{f}该如何操作呢？我们需要穿透两层函子结构。可以用\code{fmap}穿透外层的\code{Maybe}，但我们不能直接将\code{f}送入\code{Maybe}内部，因为\code{f}无法处理列表。必须将\code{(fmap f)}应用于内层列表。例如，让我们看看如何对Maybe整数列表中的元素进行平方运算：

\src{snippet30}
编译器在分析类型后会确定：对外层\code{fmap}应该使用\code{Maybe}实例的实现，对内层则使用列表函子的实现。上述代码其实可以改写为：

\src{snippet31}
但请记住\code{fmap}可以视为单参数函数：

\src{snippet32}
在此例中，\code{(fmap . fmap)}中的第二个\code{fmap}接收参数：

\src{snippet33}
并返回如下类型的函数：

\src{snippet34}
第一个\code{fmap}再接收该函数并返回：

\src{snippet35}
最终该函数作用于\code{mis}。因此两个函子的组合是一个新函子，其\code{fmap}是各分量\code{fmap}的组合。回到范畴论视角：函子组合显然具有结合律（对象映射和态射映射分别满足结合律）。每个范畴中还存在平凡的恒等函子：它将每个对象映射到自身，每个态射映射到自身。因此函子具备某个范畴中态射的所有性质。但这是哪个范畴呢？它的对象应该是范畴本身，态射应该是函子——这就是"范畴的范畴"（category of categories）。不过包含所有范畴的范畴会引发类似"所有集合的集合"那样的悖论。然而确实存在一个包含所有小范畴（small category）的范畴$\Cat$（这个范畴本身是大范畴，因此不能包含自身）。小范畴是指其对象构成一个集合（而非更大基数的类）的范畴。需要注意的是，在范畴论中即使不可数无限集也被视为"小"。我认为有必要提及这些概念，因为它揭示了抽象层级间结构重复出现的奇妙现象。后续我们将看到函子本身也能形成范畴。

\section{挑战}

\begin{enumerate}
  \tightlist
  \item
        我们能否通过定义：

        \begin{snip}{haskell}
fmap _ _ = Nothing
\end{snip}

        来将 \code{Maybe} 类型构造器转化为函子？该定义会忽略其两个参数。（提示：验证函子定律。）
  \item
        证明读者函子(reader functor)的函子定律。提示：这个问题非常简单。
  \item
        在你第二喜欢的语言中（第一选择当然是Haskell）实现读者函子。
  \item
        证明列表函子(list functor)的函子定律。假设列表尾部的函子定律成立（换句话说，使用\emph{数学归纳法}进行证明）。
\end{enumerate}